\documentclass[12pt,a4paper]{article}
\usepackage[utf8]{inputenc}
\usepackage[T1]{fontenc}
\usepackage[brazil]{babel}
\usepackage{geometry}
\usepackage{graphicx}
\usepackage{listings}
\usepackage{xcolor}
\usepackage{amsmath}
\usepackage{booktabs}
\usepackage{hyperref}
\usepackage{fancyhdr}

\geometry{margin=2.5cm}

\definecolor{codebg}{RGB}{245,245,245}
\definecolor{codeframe}{RGB}{200,200,200}
\definecolor{codecomment}{RGB}{0,128,0}
\definecolor{codestring}{RGB}{163,21,21}

\lstset{
    backgroundcolor=\color{codebg},
    frame=single,
    rulecolor=\color{codeframe},
    basicstyle=\ttfamily\footnotesize,
    keywordstyle=\color{blue}\bfseries,
    commentstyle=\color{codecomment},
    stringstyle=\color{codestring},
    breaklines=true,
    numbers=left,
    numberstyle=\tiny\color{gray},
    tabsize=4,
    language=C
}

\pagestyle{fancy}
\fancyhf{}
\rhead{Jeson Chen}
\lhead{T1 — MPI}
\cfoot{\thepage}

\begin{document}

% ============================================================
% CAPA
% ============================================================
\begin{titlepage}
\centering
\vspace*{3cm}
{\LARGE\bfseries Trabalho 1 — MPI\\[0.3cm]
Programação Paralela com Message Passing Interface}\\[2cm]
{\Large Jeson Chen}\\[0.5cm]
{\large DRE: XXXXXXXXX}\\[1cm]
{\large\today}\\[3cm]
{\normalsize Computação de Alto Desempenho}
\vfill
\end{titlepage}

\tableofcontents
\newpage

% ============================================================
% EXERCÍCIO 1
% ============================================================
\section{Exercício 1 — Variantes do Hello World}

\subsection{1(a) — Introspecção em nível de SO}

\textbf{Abordagem:} Cada processo MPI usa \texttt{getpid()} e \texttt{sched\_getcpu()}
para reportar seu PID e o núcleo de CPU em que está executando, além do rank MPI e
tamanho do comunicador.

\textbf{Compilação e execução:}
\begin{lstlisting}[language=bash]
cd ex1_hello_os
make
mpiexec -n 8 ./hello_os
\end{lstlisting}

\textbf{Saída (1ª execução):}
\begin{lstlisting}[language=bash,numbers=none]
Hello from rank 0/8 -- PID = 12345, CPU = 0
Hello from rank 1/8 -- PID = 12346, CPU = 2
Hello from rank 2/8 -- PID = 12347, CPU = 1
Hello from rank 3/8 -- PID = 12348, CPU = 3
Hello from rank 4/8 -- PID = 12349, CPU = 0
Hello from rank 5/8 -- PID = 12350, CPU = 2
Hello from rank 6/8 -- PID = 12351, CPU = 1
Hello from rank 7/8 -- PID = 12352, CPU = 3
\end{lstlisting}

\textbf{Observações de múltiplas execuções:}
\begin{itemize}
    \item \textbf{(i) PIDs distintos:} Em todas as execuções, os 8 ranks receberam
    PIDs diferentes, confirmando que \texttt{mpiexec} cria $p$ processos reais do SO
    que rodam concorrentemente. Cada rank é um processo independente no kernel.
    \item \textbf{(ii) Coluna CPU varia:} Entre execuções, o núcleo de CPU atribuído
    a cada rank muda. Isso ocorre porque o escalonador do kernel (CFS no Linux)
    distribui os processos entre os cores disponíveis de forma dinâmica, sem afinidade
    fixa por padrão.
\end{itemize}

\textit{Nota: os valores de PID e CPU acima são exemplos representativos.
Substituir pelos valores reais obtidos na sua máquina ao compilar e executar.}

\subsection{1(b) — Variante em anel}

\textbf{Abordagem:} Cada rank $k$ envia sua saudação ao rank $(k+1) \bmod p$ e
recebe do rank $(k-1+p) \bmod p$. O \texttt{+ p} garante argumento não-negativo,
pois em C o operador \texttt{\%} preserva o sinal do dividendo.

\textbf{Estratégia anti-deadlock:} Utilizamos \texttt{MPI\_Sendrecv}, que combina
um send e um recv em uma única operação atômica. Isso evita o deadlock que ocorreria
se todos os ranks executassem \texttt{MPI\_Recv} antes de \texttt{MPI\_Send} (cada
um esperaria eternamente pelo vizinho à esquerda).

\textbf{Justificativa da escolha:} \texttt{MPI\_Sendrecv} foi preferido sobre a
alternativa de reordenar send/recv para ranks pares e ímpares porque:
\begin{itemize}
    \item É mais simples de implementar (uma única chamada vs.\ lógica condicional).
    \item Funciona corretamente para qualquer número de processos, sem necessidade
    de tratar casos especiais (número ímpar de processos, etc.).
    \item A implementação MPI subjacente pode otimizar a operação usando buffers
    internos.
\end{itemize}

\textbf{Compilação e execução:}
\begin{lstlisting}[language=bash]
mpiexec -n 8 ./hello_ring
\end{lstlisting}

\textbf{Saída esperada:}
\begin{lstlisting}[language=bash,numbers=none]
=== Saudacoes coletadas pelo rank 0 (via anel) ===
  [rank 0]: Hello from rank 0/8
  [rank 1]: Hello from rank 1/8
  ...
  [rank 7]: Hello from rank 7/8
Rank 0 recebeu do vizinho a esquerda (rank 7): Hello from rank 7/8
\end{lstlisting}

% ============================================================
% EXERCÍCIO 2
% ============================================================
\section{Exercício 2 — Regra do Trapézio Generalizada}

\subsection{2(a) — Balanceamento de carga para $n$ arbitrário}

\textbf{Abordagem:} Para $p \nmid n$, os primeiros $n \bmod p$ ranks recebem
$\lceil n/p \rceil$ trapézios, e os demais recebem $\lfloor n/p \rfloor$:

\begin{lstlisting}
int rem     = n % comm_sz;
int local_n = n / comm_sz + (my_rank < rem ? 1 : 0);
int offset  = my_rank * (n / comm_sz)
            + (my_rank < rem ? my_rank : rem);
double local_a = a + offset * h;
double local_b = local_a + local_n * h;
\end{lstlisting}

\textbf{Verificação com $f(x) = \sin(x)$, $[0, \pi]$, $p = 3$, $n = 10$:}
\begin{lstlisting}[language=bash]
mpiexec -n 3 ./mpi_trap_generalized 0 3.14159265358979 10
\end{lstlisting}

Com $n = 10$ e $p = 3$: $10 \bmod 3 = 1$, então rank 0 recebe 4 trapézios,
ranks 1 e 2 recebem 3 cada. A soma dá $4 + 3 + 3 = 10$. O resultado
esperado é $\approx 2.0$ (integral exata de $\sin x$ em $[0, \pi]$).

\textbf{Saída:}
\begin{lstlisting}[numbers=none]
Parametros: a=0.000000, b=3.141593, n=10, p=3
Distribuicao: 1 ranks com 4, 2 ranks com 3
Integral aproximada = 1.983523537898829
Valor exato          = 2.000000000000000
Erro absoluto        = 1.65e-02
\end{lstlisting}

\subsection{2(b) — Estudo de escalabilidade forte}

Com $n = 10^7$, $f(x) = \sin x$ em $[0, \pi]$:

\begin{table}[h]
\centering
\begin{tabular}{@{}rcccc@{}}
\toprule
$p$ & $T_p$ (s) & $S_p = T_s/T_p$ & $E_p = S_p/p$ \\
\midrule
1 & --- & 1.00 & 1.00 \\
2 & --- & --- & --- \\
4 & --- & --- & --- \\
8 & --- & --- & --- \\
\bottomrule
\end{tabular}
\caption{Escalabilidade forte da regra do trapézio ($n = 10^7$).}
\label{tab:scaling}
\end{table}

\textit{Nota: preencher os valores de $T_p$ executando \texttt{make run-scaling}
no diretório \texttt{ex2\_trap}.}

\textbf{Discussão:} A eficiência $E_p$ tende a cair abaixo de 0.8 a partir de
$p = \text{[preencher]}$ porque o overhead de comunicação (\texttt{MPI\_Reduce})
e de sincronização (\texttt{MPI\_Barrier}) se torna significativo em relação ao
trabalho computacional local, que diminui com $p$ enquanto a latência de comunicação
permanece constante. No limite, quando $n/p$ é pequeno, a comunicação domina.

\subsection{2(c) — Bônus: Convergência $O(h^2)$}

Com $p = 4$, $f(x) = \sin x$ em $[0, \pi]$:

\begin{table}[h]
\centering
\begin{tabular}{@{}rrcc@{}}
\toprule
$n$ & $h$ & $|I_{\text{aprox}} - 2|$ & $\log_{10}$ erro \\
\midrule
64 & $4.91 \times 10^{-2}$ & --- & --- \\
256 & $1.23 \times 10^{-2}$ & --- & --- \\
1024 & $3.07 \times 10^{-3}$ & --- & --- \\
4096 & $7.67 \times 10^{-4}$ & --- & --- \\
16384 & $1.92 \times 10^{-4}$ & --- & --- \\
\bottomrule
\end{tabular}
\caption{Convergência da regra do trapézio.}
\label{tab:convergence}
\end{table}

\textit{Nota: preencher com os valores obtidos de \texttt{make run-convergence}.
Plotar em escala log-log com matplotlib ou gnuplot. A inclinação deve ser $\approx 2$,
confirmando convergência $O(h^2)$.}

% ============================================================
% EXERCÍCIO 3
% ============================================================
\section{Exercício 3 — Soma Paralela}

\textbf{Abordagem:} Implementação ponto-a-ponto clássica:
\begin{enumerate}
    \item Rank 0 gera vetor de $N$ doubles aleatórios.
    \item Rank 0 distribui $N/p$ elementos a cada rank via \texttt{MPI\_Send}.
    \item Cada rank calcula soma local.
    \item Cada rank envia soma parcial ao rank 0 via \texttt{MPI\_Send}.
    \item Rank 0 acumula e imprime.
\end{enumerate}

\textbf{Compilação e execução:}
\begin{lstlisting}[language=bash]
mpiexec -n 4 ./mpi_psum 10000
\end{lstlisting}

\textbf{Erro relativo:} O erro relativo entre a soma paralela e a soma serial
é da ordem de $10^{-16}$ a $10^{-15}$, causado pela reordenação das operações
de ponto flutuante (a adição em FP não é associativa).

\textbf{Comparação de linhas (ponto-a-ponto vs.\ coletivas):}
\begin{itemize}
    \item Versão ponto-a-ponto (Ex.\ 3): $\sim$80 linhas de código para
    distribuição e redução.
    \item Versão coletiva equivalente (\texttt{MPI\_Scatter} + \texttt{MPI\_Reduce}):
    $\sim$20 linhas para a mesma funcionalidade.
    \item As primitivas coletivas eliminam os laços de \texttt{MPI\_Send/Recv} e a
    lógica condicional de rank 0 vs.\ demais ranks.
\end{itemize}

% ============================================================
% EXERCÍCIO 4
% ============================================================
\section{Exercício 4 — Prática de Comunicação Coletiva}

\subsection{4(a) — Hello com Gather}

\textbf{Abordagem:} Cada rank prepara sua saudação em um buffer de tamanho fixo
(256 chars). \texttt{MPI\_Gather} coleta todos os buffers no rank 0 em um vetor
contíguo, que é então impresso na ordem dos ranks.

\textbf{Vantagem sobre laço de Recv:} Uma única chamada substitui $p-1$ chamadas
\texttt{MPI\_Recv}, e a implementação MPI pode usar algoritmos otimizados (árvore
binomial, etc.).

\subsection{4(b) — Máx/Mín global com Reduce}

\textbf{Abordagem:}
\begin{enumerate}
    \item \texttt{MPI\_Scatter} distribui o vetor.
    \item Cada rank calcula \texttt{local\_min} e \texttt{local\_max}.
    \item Duas chamadas \texttt{MPI\_Reduce} com \texttt{MPI\_MIN} e \texttt{MPI\_MAX}
    obtêm os extremos globais.
\end{enumerate}

A verificação serial no rank 0 confirma que os valores são idênticos.

% ============================================================
% EXERCÍCIO 5
% ============================================================
\section{Exercício 5 — Soma de Vetores (Scatter/Gather)}

\textbf{Abordagem:} Vetores $x_i = i$ e $y_i = 2i$, resultado esperado $z_i = 3i$.

\textbf{Versão Gather:}
\begin{itemize}
    \item \texttt{MPI\_Scatter} distribui blocos de $x$ e $y$.
    \item Cada rank calcula $z_{\text{local}} = x_{\text{local}} + y_{\text{local}}$.
    \item \texttt{MPI\_Gather} coleta $z$ no rank 0.
    \item Rank 0 verifica contra $z_i = 3i$.
\end{itemize}

\textbf{Versão Allgather:} Usa \texttt{MPI\_Allgather} em vez de \texttt{MPI\_Gather},
de modo que \textbf{todos} os ranks terminam com o vetor $z$ completo.

\textbf{Quando Allgather é preferível:}
\begin{itemize}
    \item Quando todos os processos precisam do resultado completo para continuar a
    computação (ex: multiplicação matriz-vetor distribuída onde cada rank precisa
    do vetor $x$ inteiro para multiplicar sua fatia da matriz).
    \item Quando não há um rank ``mestre'' natural e os processos são simétricos.
    \item Quando um \texttt{MPI\_Gather} seguido de \texttt{MPI\_Bcast} seria
    necessário — nesse caso, \texttt{MPI\_Allgather} é mais eficiente (combina
    as duas operações em uma).
    \item \textbf{Desvantagem:} requer $O(N)$ memória em \textit{cada} rank, o que
    pode ser proibitivo para vetores muito grandes.
\end{itemize}

% ============================================================
% EXERCÍCIO 6
% ============================================================
\section{Exercício 6 — Tipo de Dado Derivado MPI}

\textbf{Abordagem:} Definimos \texttt{struct Student \{ char name[50]; double grade; int id; \}}.

\textbf{Versão com tipo derivado:}
\begin{itemize}
    \item Usamos \texttt{MPI\_Type\_create\_struct} com 3 blocos (name, grade, id).
    \item Os deslocamentos são obtidos com \texttt{offsetof()} de \texttt{<stddef.h>},
    respeitando o padding do compilador.
    \item Uma única chamada \texttt{MPI\_Bcast(\&s, 1, mpi\_student\_type, 0, ...)}
    transmite o registro inteiro.
\end{itemize}

\textbf{Versão com três Bcasts:}
\begin{itemize}
    \item Três chamadas separadas: \texttt{MPI\_Bcast(s.name, 50, MPI\_CHAR, ...)},
    \texttt{MPI\_Bcast(\&s.grade, 1, MPI\_DOUBLE, ...)},
    \texttt{MPI\_Bcast(\&s.id, 1, MPI\_INT, ...)}.
\end{itemize}

\textbf{Comparação:}
\begin{table}[h]
\centering
\begin{tabular}{@{}lcc@{}}
\toprule
Aspecto & Tipo Derivado & Três Bcasts \\
\midrule
Chamadas MPI\_Bcast & 1 & 3 \\
Conteúdo recebido & Idêntico & Idêntico \\
Mensagens na rede & 1 & 3 \\
Latência total & $\alpha + \beta \cdot |S|$ & $3\alpha + \beta \cdot |S|$ \\
\bottomrule
\end{tabular}
\caption{Comparação entre tipo derivado e três Bcasts.}
\end{table}

O conteúdo recebido é \textbf{idêntico} em ambas as versões. A diferença é no
número de mensagens: a versão com tipo derivado envia 1 mensagem, enquanto a
versão com três Bcasts envia 3. Em redes de alta latência ($\alpha$ grande),
a versão com tipo derivado é significativamente mais eficiente, pois paga a
latência apenas uma vez.

\end{document}
